% !TEX root =  Lecture19.tex


%%%%%%%%%%%%%%%%%%%%%%%%%%%%%%%%%%%%
% Recap/Agenda
%%%%%%%%%%%%%%%%%%%%%%%%%%%%%%%%%%%%
\begin{frame}
\frametitle{Module 4: Bayesian Statistics}
    \begin{itemize}
        \item \hl{Previously:} conjugate models: the posterior in closed form, its mean a weighted average of prior and data.
        \item \hl{This time:} what to do when there is \emph{no} closed form (the usual case), so we must \hl{approximate} the posterior.
          \begin{itemize}
              \item why the normalizing constant is the obstacle;
              \item \hl{grid approximation}, turning the integral into a sum;
              \item \hl{Markov chain Monte Carlo}, sampling from the posterior instead;
              \item reading the output: \hl{trace plots}, burn-in, step size, diagnostics.
          \end{itemize}
        \item \hl{Reading:} Bayes Rules!\ Chapters 6 and 7.
        \item \hl{Homework:}
    \end{itemize}
\end{frame}


%%%%%%%%%%%%%%%%%%%%%%%%%%%%%%%%%%%%
% Learning objectives:
%%%%%%%%%%%%%%%%%%%%%%%%%%%%%%%%%%%%
\begin{frame}
    \frametitle{Lecture Objectives}
    \goals{%
    \begin{itemize}
    \item explain why a non-conjugate model has \hl{no closed-form posterior}, and identify the normalizing constant as the obstacle;
    \item use \hl{grid approximation} to turn the posterior integral into a sum, and say why it breaks down as the number of parameters grows;
    \item describe how \hl{Metropolis--Hastings} MCMC draws posterior samples via propose / accept-or-reject, and read an \hl{acceptance probability};
    \item diagnose a chain from its \hl{trace plot}: estimate the \hl{burn-in}, spot a \hl{step size} that is too large or too small;
    \item say what \hl{R-hat}, \hl{ESS}, and \hl{MCSE} each tell you about whether to trust the output.
    \end{itemize}}
\end{frame}


%%%%%%%%%%%%%%%%%%%%%%%%%%%%%%%%%%%%
\begin{frame}
    \frametitle{Lecture Objectives}
    \objectives{%
    \begin{itemize}
    \item \textbf{(M4, LO5)} Using Markov chain Monte Carlo
    \end{itemize}}
    \vspace{0.3em}
    {\small Building on, and revisiting:}
    \objectives{%
    \begin{itemize}
    \item \textbf{(M4, LO3)} Conjugate Models: the whale-sightings posterior from Lecture 18, now recomputed \emph{without} conjugacy
    \end{itemize}}
\end{frame}
